\section{Tutorial 9: Gradient Flow Experiments}

\subsection{Objectives}

This tutorial provides hands-on and conceptual understanding of gradient flow in deep networks:

\begin{itemize}
    \item Observe vanishing and exploding gradients
    \item Compare activation functions
    \item Study the effect of initialization
    \item Understand the role of normalization
\end{itemize}

\medskip

\noindent
\textbf{Focus.}  
Turn theoretical insights about gradient flow into observable behavior.

\subsection{Exercise 1: Gradient Norm Across Depth}

Consider a deep neural network with $L$ layers.

\medskip

\noindent
\textbf{Task.}
\begin{enumerate}
    \item For each layer $\ell$, define:
    \[
    g_\ell = \left\|\frac{\partial L}{\partial h_\ell}\right\|.
    \]
    \item Predict how $g_\ell$ behaves as $\ell$ decreases (toward input).
\end{enumerate}

\medskip

\noindent
\textbf{Discussion.}

\begin{itemize}
    \item If gradients vanish:
    \[
    g_1 \ll g_L
    \]

    \item If gradients explode:
    \[
    g_1 \gg g_L
    \]
\end{itemize}

\medskip

\noindent
\textbf{Insight.}  
Gradient norms reveal how learning signals propagate.

\subsection{Exercise 2: Activation Function Comparison}

Compare sigmoid and ReLU networks of depth $L$.

\medskip

\noindent
\textbf{Tasks.}
\begin{enumerate}
    \item Explain why sigmoid leads to vanishing gradients.
    \item Explain why ReLU improves gradient flow.
\end{enumerate}

\medskip

\noindent
\textbf{Solution.}

\begin{itemize}
    \item Sigmoid:
    \begin{itemize}
        \item $\sigma'(z) \leq 0.25$
        \item repeated multiplication $\Rightarrow$ exponential decay
    \end{itemize}

    \item ReLU:
    \begin{itemize}
        \item derivative is $1$ in active region
        \item preserves gradient magnitude
    \end{itemize}
\end{itemize}

\medskip

\noindent
\textbf{Insight.}  
Activation choice directly affects gradient dynamics.

\subsection{Exercise 3: Initialization Effects}

Assume weights are initialized as:
\[
W_{ij} \sim \mathcal{N}(0, \sigma^2).
\]

\medskip

\noindent
\textbf{Tasks.}
\begin{enumerate}
    \item What happens if $\sigma$ is too small?
    \item What happens if $\sigma$ is too large?
\end{enumerate}

\medskip

\noindent
\textbf{Solution.}

\begin{itemize}
    \item Small $\sigma$:
    \begin{itemize}
        \item activations shrink
        \item gradients vanish
    \end{itemize}

    \item Large $\sigma$:
    \begin{itemize}
        \item activations grow
        \item gradients explode
    \end{itemize}
\end{itemize}

\medskip

\noindent
\textbf{Insight.}  
Initialization controls both forward and backward signal propagation.

\subsection{Exercise 4: Effect of Depth}

Consider increasing depth $L$.

\medskip

\noindent
\textbf{Tasks.}
\begin{enumerate}
    \item Describe how gradient stability changes with depth.
    \item Why are shallow networks easier to train?
\end{enumerate}

\medskip

\noindent
\textbf{Solution.}

\begin{itemize}
    \item Gradient instability increases with depth
    \item Fewer multiplicative terms in shallow networks
\end{itemize}

\medskip

\noindent
\textbf{Insight.}  
Depth amplifies both expressivity and instability.

\subsection{Exercise 5: Normalization Effects}

Consider inserting batch normalization after each layer.

\medskip

\noindent
\textbf{Tasks.}
\begin{enumerate}
    \item How does normalization affect activations?
    \item How does it affect gradients?
\end{enumerate}

\medskip

\noindent
\textbf{Solution.}

\begin{itemize}
    \item Activations are rescaled to stable distributions
    \item Gradients become more stable and less sensitive to scale
\end{itemize}

\medskip

\noindent
\textbf{Insight.}  
Normalization stabilizes both forward and backward passes.

\subsection{Exercise 6: Practical Experiment (Conceptual)}

\textbf{Experiment design.}

Train three networks with identical architecture:

\begin{itemize}
    \item Model A: sigmoid activation
    \item Model B: ReLU activation
    \item Model C: ReLU + batch normalization
\end{itemize}

\medskip

\noindent
\textbf{Tasks.}
\begin{enumerate}
    \item Predict training behavior of each model.
    \item Rank them by training speed and stability.
\end{enumerate}

\medskip

\noindent
\textbf{Expected outcome.}

\begin{itemize}
    \item Model A:
    \begin{itemize}
        \item slow training, vanishing gradients
    \end{itemize}

    \item Model B:
    \begin{itemize}
        \item faster, but may still be unstable
    \end{itemize}

    \item Model C:
    \begin{itemize}
        \item fastest and most stable
    \end{itemize}
\end{itemize}

\medskip

\noindent
\textbf{Insight.}  
Modern deep learning relies on combining multiple stabilization techniques.

\subsection{Exercise 7 (Discussion): Diagnosing Training Failures}

\textbf{Question.}  
How can you diagnose gradient-related issues in practice?

\medskip

\noindent
\textbf{Discussion points.}

\begin{itemize}
    \item monitor gradient norms per layer
    \item observe training loss behavior
    \item check activation distributions
\end{itemize}

\medskip

\noindent
\textbf{Insight.}  
Understanding gradient flow is essential for debugging deep models.

\subsection{Summary}

\begin{itemize}
    \item Gradient flow determines whether deep networks can learn
    \item Activation functions, initialization, and normalization all affect it
    \item Instability grows with depth
    \item Practical success depends on controlling these dynamics
\end{itemize}

\medskip

\noindent
\textbf{Preparation for next tutorial.}  
We next study how architectural changes (e.g., residual connections) further improve stability and enable very deep networks.